%=======================================================================
%  CSE558 Data Science - Midterm Project Report
%  Monsoon 2026, IIIT-Delhi
%
%  Built on the course template (CSE558_Project_Report_Template.tex).
%  Figures come from ../figures/*.pdf, written by run_all.py at the exact
%  column width (83.5 mm), so they are not rescaled here.
%  Text still in LIGHT BLUE is a placeholder that must be replaced.
%  Compile with pdfLaTeX (Overleaf's default setting works), twice.
%=======================================================================
\documentclass[10pt,a4paper,twocolumn]{article}

%========================= EDIT THIS =========================
% Your project title: it sets the heading AND the PDF's title.
\newcommand{\projecttitle}{Do pre-war buildings still sell at a premium?
  Evidence from 845,607 New York City property sales}
\newif\ifshowhints
\showhintsfalse     % \showhintstrue shows every grey instruction line
%=============================================================

%-------------------------------------------- fonts (no need to edit)
\usepackage{iftex}
\ifPDFTeX
  \usepackage[T1]{fontenc}
  \usepackage[utf8]{inputenc}
  \IfFileExists{carlito.sty}%
    {\usepackage[sfdefault]{carlito}}%              Calibri look
    {\usepackage[scaled=0.92]{helvet}\renewcommand{\familydefault}{\sfdefault}}%
  \usepackage{courier}
\else
  \usepackage[no-math]{fontspec}
  \setmainfont{Carlito}[Extension=.ttf, UprightFont=*-Regular, BoldFont=*-Bold,
    ItalicFont=*-Italic, BoldItalicFont=*-BoldItalic]
\fi

%----------------------------------------------------------- packages
\usepackage[margin=18mm,bottom=20mm,columnsep=7mm]{geometry}
\usepackage{amsmath}
\usepackage{xcolor}
\usepackage{graphicx}
\usepackage{tabularx}
\usepackage{array}
\usepackage{booktabs}
\usepackage{enumitem}
\usepackage{titlesec}
\usepackage{needspace}
\usepackage[font=small,labelfont=bf]{caption}
\usepackage[hidelinks]{hyperref}
\hypersetup{pdftitle={\projecttitle},
            pdfsubject={CSE558 Data Science, Monsoon 2026: Midterm Project Report}}

%------------------------------------------------------------- layout
\definecolor{eg}{HTML}{7396C4}         % worked example (light blue)
\definecolor{hintgrey}{gray}{0.58}     % instructions (light grey)

\titleformat{\section}{\normalfont\large\bfseries}{\thesection.}{0.5em}{}
\titlespacing*{\section}{0pt}{12pt}{3pt}
\titleformat{\subsection}{\normalfont\normalsize\bfseries}{\thesubsection}{0.5em}{}
\titlespacing*{\subsection}{0pt}{10pt}{2pt}
\newcommand{\sectionbreak}{\needspace{10\baselineskip}}
\newcommand{\subsectionbreak}{\needspace{10\baselineskip}}

\setlength{\parindent}{0pt}
\setlength{\parskip}{4pt}
\setlist{leftmargin=1.2em, itemsep=2pt, topsep=2pt, parsep=0pt}
\renewcommand{\arraystretch}{1.25}
\setlength{\tabcolsep}{4pt}
\newcolumntype{L}[1]{>{\raggedright\arraybackslash}p{#1}}
\newcolumntype{Y}{>{\raggedright\arraybackslash}X}
\newcolumntype{K}[1]{>{\raggedright\arraybackslash\bfseries\boldmath}p{#1}}  % bold labels

%--------------------------------------------------------- shorthands
% Instruction line. Delete these as you fill the template in.
\newcommand{\hint}[1]{\ifshowhints{\footnotesize\itshape\color{hintgrey}#1\par}\fi}
% Worked example: inline or in a table cell. REPLACE the contents.
\newcommand{\eg}[1]{\textcolor{eg}{#1}}
% Worked example running over a whole paragraph.
\newcommand{\egblock}[1]{{\color{eg}#1\par}}
% Bold label followed by an example.
\newcommand{\field}[2]{\textbf{#1}~\eg{#2}\par}
% Part a)-c) of a test, and a small bold heading.
\newcommand{\tpart}[2]{\par\needspace{12\baselineskip}\smallskip\textbf{#1)~#2}\par}
\newcommand{\minihead}[1]{\par\needspace{8\baselineskip}\smallskip\textbf{#1}\par}
% Column and file names.
\newcommand{\code}[1]{\texttt{#1}}
% \figfile{file}{height}: your plot, or a labelled box until the file exists.
\newcommand{\figfile}[2]{%
  \IfFileExists{#1}%
    {\includegraphics[width=\linewidth,height=#2,keepaspectratio]{#1}}%
    {\fbox{\parbox[c][#2][c]{\dimexpr\linewidth-2\fboxsep-2\fboxrule\relax}%
      {\centering\small\itshape\color{hintgrey}Your plot: \detokenize{#1}}}}}
% Our own real-content label (black, unlike \field which is a placeholder).
\newcommand{\lab}[2]{\textbf{#1}~#2\par}
% Figures live in ../figures/ in the repo, or in figures/ next to main.tex on Overleaf.
\newcommand{\figdir}{../figures/}
\IfFileExists{figures/f00_analysis_funnel.pdf}{\renewcommand{\figdir}{figures/}}{}

%=======================================================================
\begin{document}

\twocolumn[%
{\small\bfseries CSE558 Data Science\ \ \textbullet\ \ Midterm Project Report\ \ \textbullet\ \ Monsoon 2026\par}
\vspace{4pt}
{\LARGE\bfseries\raggedright\projecttitle\par}
\vspace{6pt}
\rule{\linewidth}{1pt}

{\small
\begin{tabularx}{\linewidth}{@{}K{2.4cm} Y@{}}
\toprule
Group name & \eg{TODO group name} \\
Dataset & \href{https://data.cityofnewyork.us/City-Government/NYC-Citywide-Annualized-Calendar-Sales-Update/w2pb-icbu}{data.cityofnewyork.us, dataset w2pb-icbu} (NYC Department of Finance, \emph{NYC Citywide Annualized Calendar Sales Update}) \\
Code & \href{https://github.com/ami4go/Data_Science_Project_NYC_Property_Sales}{github.com/ami4go/Data\_Science\_Project\_NYC\_Property\_Sales} (scripts, README, frozen data link) \\
\bottomrule
\end{tabularx}}

\vspace{2pt}
{\small
\begin{tabularx}{\linewidth}{@{}L{0.5cm} L{3.4cm} L{2.3cm} L{4.8cm} Y@{}}
\toprule
\textbf{\#} & \textbf{Name} & \textbf{Roll no.} & \textbf{Institute email} & \textbf{Contribution} \\
\midrule
1 & \eg{TODO} & \eg{TODO} & \eg{TODO} & Dataset, raw audit and size check (1) \\
2 & \eg{TODO} & \eg{TODO} & \eg{TODO} & Cleaning, multi-lot rule and the 70/30 split; before and after (2, 3) \\
3 & \eg{TODO} & \eg{TODO} & \eg{TODO} & EDA figures and problem statements (4) \\
4 & \eg{TODO} & \eg{TODO} & \eg{TODO} & Test plan; Tests 1--2: pre-war premium and mansion tax (5) \\
5 & \eg{TODO} & \eg{TODO} & \eg{TODO} & Tests 3--4: transfers and mixed use; limitations, report and README (5, 6) \\
\bottomrule
\end{tabularx}}
\vspace{12pt}
]

%-----------------------------------------------------------------------
\section{Dataset}

\lab{Source:}{NYC Department of Finance, \emph{NYC Citywide Annualized Calendar Sales Update},
NYC Open Data (dataset \code{w2pb-icbu}), downloaded on 20 October 2026 and frozen, because the
city keeps updating it (SHA-256 \code{cf8d76c1\ldots}; the code refuses any other file). One row
is one recorded sale: a deed on one tax lot, dated 1 Jan 2016 to 31 Dec 2025.}

\lab{What is raw about it:}{31.0\% of the ``sales'' are \$0 ownership transfers; floor and land
area are stored as text with commas (\code{"55,000"}) and are 0 or blank in 44.3\% of rows; one
building category is spelled two ways (91 labels for 47 types); the borough is a code from 1 to 5;
the \code{EASE-MENT} column is entirely blank; 3,386 rows are exact duplicates; and a deal covering
several lots repeats its \emph{total} price on every lot. Prices run from \$0 to \$2.4 billion.}

{\small
\begin{tabularx}{\linewidth}{@{}L{3.0cm} Y Y@{}}
\toprule
 & \textbf{Raw file} & \textbf{After encoding} \\
\midrule
Instances (rows) & 845,607 & 566,956 market sales \\
Attributes (columns) & 29 & 71 \\
Instances $\times$ attributes & 24,522,603 & 40,253,876 \\
\bottomrule
\end{tabularx}}

\lab{Size check:}{$566{,}956\times71=40{,}253{,}876\ge3{,}000{,}000$ (13.4 times the minimum).
The raw table already clears 3M; one-hot encoding borough (5 levels), building category (47) and
tax class (4) adds 53 columns to 15 numeric and derived ones. Even our smallest test population,
the 248,409 sales behind P1, gives $248{,}409\times71=17.6$M.}

\minihead{Key features}
{\footnotesize
\begin{tabularx}{\linewidth}{@{}L{2.6cm} L{1.4cm} Y@{}}
\toprule
\textbf{Feature} & \textbf{Type} & \textbf{Real-world meaning} \\
\midrule
\code{sale\_price} & continuous & recorded price, US\$; response for P2--P4 \\
\code{price\_per\_sqft} & continuous & derived: price / gross floor area; response for P1 \\
\code{gross\_sqft} & continuous & floor area of the building, sq ft \\
\code{year\_built} & discrete & assessor's year built; pre-war = before 1940 \\
\code{total\_units} & discrete & residential + commercial units on the lot \\
\code{borough} & nominal & 5 boroughs \\
\code{neighborhood} & nominal & 262 assessor neighbourhoods \\
\code{building\_category} & nominal & 47 types, e.g.\ one-family, condo elevator apartment \\
\code{latitude}, \code{longitude} & continuous & location; used for the toll-zone flag \\
\code{sale\_date} & date & 2016-01-01 to 2025-12-31 \\
\bottomrule
\end{tabularx}}

%-----------------------------------------------------------------------
% The template's \needspace before sections over-reserves here and would leave
% page 1 half empty, so it is switched off for this one heading.
{\let\sectionbreak\relax\section{Data preparation}}

{\footnotesize
\begin{tabularx}{\linewidth}{@{}L{1.95cm} L{2.2cm} Y@{}}
\toprule
\textbf{Issue found} & \textbf{What we did} & \textbf{Why, and not the alternative} \\
\midrule
\$0--\$100 sales (274,201; 32.4\%)
 & flagged \emph{nominal}; kept only for the transfer test (T3)
 & transfers between relatives or companies are not prices; deleting them would make T3 impossible \\
Area stored as text (\code{"55,000"})
 & commas removed, converted to numbers
 & a format fix; nothing is lost \\
Floor area 0 or blank (374,853; 44.3\%)
 & 0 set to missing; \textbf{not imputed}
 & missing by design: 95\% of condo/co-op unit sales lack it vs 5\% of others (Fig.~\ref{fig:missing}); imputing would invent P1's denominator \\
\bottomrule
\end{tabularx}}

{\footnotesize
\begin{tabularx}{\linewidth}{@{}L{1.95cm} L{2.2cm} Y@{}}
\toprule
\textbf{Issue found} & \textbf{What we did} & \textbf{Why, and not the alternative} \\
\midrule
91 category labels for 47 types
 & spaces collapsed, upper-cased
 & otherwise one-hot encoding splits one type into two columns \\
3,386 exact duplicate rows
 & dropped
 & a sale counted twice inflates counts and shrinks p-values \\
Multi-lot deals (35,044 rows): one deed repeats its total price on every lot
 & flagged; no price per sq ft; price statistics use one row per deal (566,956 rows $=$ 545,669 deals)
 & a null model separates real deals from same-day coincidences (Fig.~\ref{fig:nullmodel}); counting lots instead of deals flips the toll-zone comparison (Fig.~\ref{fig:toll}) \\
Year built 0 or before 1800 (63,355)
 & set to missing; excluded from P1 only
 & placeholders and typos; the sales are valid for other questions \\
Year built heaped on round years (Fig.~\ref{fig:heaping})
 & kept; T1 is repeated without 1935--1945
 & dropping rounded years would remove most pre-war stock \\
Price per sq ft outliers (3,833)
 & flagged outside $3\times$IQR fences on log \$/sq ft per building type, \textbf{fitted on the exploration set only}
 & per type because levels differ; fitted like a scaler, so the confirmation set is not looked at \\
Prices above \$100M (1,083)
 & kept; analysed on a log scale
 & real trophy and portfolio deals; trimming them deletes real sales \\
Toll zone (Manhattan, 60th St and below) is not a column
 & flag from coordinates, on a line fitted to the rotated street grid
 & all 24,072 sales on East/West 50th--70th St land on the correct side \\
Nominal columns
 & one-hot: borough, building category, tax class
 & label codes would impose a fake order \\
\bottomrule
\end{tabularx}}

\begin{figure}[!ht]
  \figfile{\figdir f00b_multi_lot_null_model.pdf}{5.3cm}
  \caption{Same-block lot groups sharing a day and price are real multi-lot deals (12,473 observed
  vs 543 when sale dates are shuffled); two-lot matches across blocks are coincidence (24,557 vs
  26,611), so they are not flagged.}
  \label{fig:nullmodel}
\end{figure}

%-----------------------------------------------------------------------
\section{What changed after preparation}

{\footnotesize
\begin{tabularx}{\linewidth}{@{}L{3.3cm} L{1.4cm} Y@{}}
\toprule
\textbf{Measure} & \textbf{Before} & \textbf{After} \\
\midrule
Rows & 845,607 & 842,221 \\
Attributes & 29 & 55 (28 kept + 27 derived or flags) \\
Exact duplicate rows & 3,386 & 0 \\
Building-category labels & 91 & 47 \\
Numbers stored as text & 2 columns & 0 \\
Hidden missing (0 meaning ``not recorded'') & 314,567 & 0 (now explicit) \\
Market-price rows & 571,406 (\textgreater\ \$100) & 566,956 (\textgreater\ \$10,000) $=$ 545,669 deals \\
Rows with a valid price per sq ft & -- & 248,627 \\
\bottomrule
\end{tabularx}}

Missing cells \emph{rise} from 1,793,579 to 2,097,690 (excluding the dropped blank column),
because 314,567 zeros that meant ``not recorded'' are now marked missing; left as 0 they would
bias every mean. We lost only the 3,386 duplicate rows outright, but each question can use fewer:
price analyses drop the 32.4\% transfers, and price per sq ft is defined for 29.4\% of raw rows
(Fig.~\ref{fig:funnel}), so P1 speaks for houses and rental buildings, not apartments.

\begin{figure}[!ht]
  \figfile{\figdir f00_analysis_funnel.pdf}{5.8cm}
  \caption{Only 29\% of raw rows can answer a price-per-sq-ft question: transfers remove a third,
  and missing floor area another 35\%.}
  \label{fig:funnel}
\end{figure}

%-----------------------------------------------------------------------
\section{Basic questions and problem statement}

\lab{Split before exploring:}{70\% exploration / 30\% confirmation, \textbf{seed 42}, stratified
by borough $\times$ sale year (50 strata; exploration share 0.698--0.710 in each), with all lots of
one deal in the same half: 590,043 and 252,178 rows. All figures in this section use the
exploration set; the tests in Section~\ref{sec:tests} use only the confirmation set.}

\begin{figure}[!ht]
  \figfile{\figdir f03_missingness.pdf}{9.4cm}
  \caption{Floor area is missing for 44\% of sales, but for 95\% of condo/co-op unit sales and only
  5\% of the rest: missing by design, not at random, so we do not impute it.}
  \label{fig:missing}
\end{figure}

\begin{figure}[!ht]
  \figfile{\figdir f04_year_built_heaping.pdf}{5.0cm}
  \caption{61\% of parcels have a year built ending in 0 or 5 (20\% expected); 1940 holds 12,337
  parcels vs 540 in 1939, so the pre-war cutoff needs a sensitivity check.}
  \label{fig:heaping}
\end{figure}

\begin{figure}[!ht]
  \figfile{\figdir f06_prewar_vs_newer.pdf}{9.2cm}
  \caption{Pooled, pre-war residential buildings sell for 6.6\% more per sq ft; within the same
  neighbourhood and type the median premium is only +1.6\% (383 cells, IQR $-6$\% to $+11$\%), near
  the 5\% that would matter.}
  \label{fig:prewar}
\end{figure}

\begin{figure}[!ht]
  \figfile{\figdir f10_toll_zone.pdf}{6.6cm}
  \caption{Counting each deal once, 2024--25 median prices rose 1.6\% inside the toll zone vs
  6.6\% outside; counting every lot of multi-lot deals would show $+10.4$\% vs $+5.1$\%, driven by
  one \$170M deal repeated on 280 lots.}
  \label{fig:toll}
\end{figure}

\minihead{Basic questions}
\begin{enumerate}[label=\arabic*.]
  \item Is Manhattan the most expensive borough? By median deal price yes (\$1.2M vs
  \$514K--\$885K), though the price ranges overlap heavily.
  \item Do bigger homes cost more? Yes, moderately: Spearman $\rho=0.49$ for 1--3 family homes.
  \item Have prices risen since 2016? Yes, in every borough: +22\% in Manhattan, +48--71\% elsewhere.
  \item Which neighbourhood is priciest per sq ft? Greenwich Village-West, \$2,020.
\end{enumerate}

\minihead{Our problem statements}
\noindent\fbox{\begin{minipage}{\dimexpr\linewidth-2\fboxsep-2\fboxrule\relax}
\textbf{P1 (primary).} Within the same neighbourhood and building type, do pre-war buildings
(built before 1940) sell for more per square foot than newer ones?

\smallskip
\textbf{P2.} After the July 2019 mansion-tax step at \$2M, do more residential deals land just
below \$2M than before?

\smallskip
\textbf{P3.} Holding building type fixed, does the share of \$0--\$100 ownership transfers differ
between boroughs?

\smallskip
\textbf{P4.} Within the same neighbourhood and building type, do buildings with shops sell at a
different price per square foot from purely residential ones?
\end{minipage}}

\lab{How they are novel:}{the basic questions read one column against price, and with 250K rows
any such gap is ``significant''. Ours compare like with like: the same neighbourhood and type (P1,
P4), the same building mix (P3), or the same narrow price window before and after a tax change,
with placebo windows where the tax did not change (P2). Each needs columns no raw row records
(price per sq ft, a pre-war flag, deal-level prices, a price class, a mixed-use flag).}

\lab{Correlation is not causation:}{P2 has a causal flavour: the tax pushed prices under \$2M. Two
things support that reading: the comparison is with the same window before the tax, and windows
at \$1.5M and \$2.5M, where the tax did not step, act as placebos. It still cannot tell a real price
cut from money moved off the deed (furniture, closing credits). P1, P3 and P4 are associations:
pre-war buildings may also differ in renovation, and mixed-use buildings sit on busier streets.}

\lab{Hypotheses we dropped:}{COVID and Manhattan prices (too few Manhattan per-sq-ft sales), the
congestion toll (no stable pre-toll trend; Section~6), new buildings (too few sales), and pooled
questions such as tax class by borough or older vs newer citywide, which any test would reject.}

%-----------------------------------------------------------------------
\section{Statistical inferences}\label{sec:tests}

\lab{Rows tested:}{only the confirmation set (252,178 rows), with every hypothesis, test and
threshold fixed in \code{reports/test\_plan.md} before the tests were run. All tests are from the
course (z, t, paired t, F and $\chi^2$); $\alpha=0.05$. Fig.~\ref{fig:tests} shows every estimate
against the size that would matter.}

\begin{figure}[!ht]
  \figfile{\figdir f13_test_results.pdf}{10.2cm}
  \caption{Estimates with 95\% CIs; shaded = smaller than the size that matters. Only T2 at \$2M,
  T4, and Bronx / Staten Island in T3 clear it; the pre-war premium (T1) does not.}
  \label{fig:tests}
\end{figure}

%------------------------------------------------------------ Test 1
\subsection*{Test 1 (P1): do pre-war buildings sell for more per sq ft, like for like?}

\tpart{a}{Hypotheses, written before running the test}
{\small
\begin{tabularx}{\linewidth}{@{}K{3.4cm} Y@{}}
\toprule
$H_0$ & $\mu_d=0$ \\
$H_1$ & $\mu_d>0$ \\
Significance level & $\alpha=0.05$ \\
Side & one-sided (right) \\
Test & paired t-test \\
Smallest difference that matters & 5\% ($\mu_d=\ln 1.05=0.049$) \\
\bottomrule
\end{tabularx}}

$d$ = mean log \$/sq ft of pre-war minus newer sales in one neighbourhood $\times$ building-type
cell (1--3 family and rental buildings; cells with $\ge5$ of each). One-sided, because P1 claims a
premium. 5\%, because a smaller gap is lost in normal price negotiation.

\tpart{b}{Other tests we considered}
Each cell is a matched pair (same place, same type), $\sigma$ is unknown, and 367 cells make the
mean difference near normal: a paired t-test. Pairing removes location, which drives most of the
raw gap (Fig.~\ref{fig:prewar}).
{\footnotesize
\begin{tabularx}{\linewidth}{@{}L{2.4cm} Y@{}}
\toprule
\textbf{Also considered} & \textbf{Why we did not use it} \\
\midrule
Welch two-sample t on all sales & ignores neighbourhood, so it measures location mix. Run anyway: +7.1\% (6.3\% to 8.0\%), $t=16.75$, d.f.\ $=71{,}025$, $p<10^{-60}$. \\
Pooled two-sample t & needs equal variances; the F-test rejects them ($F=1.60$, d.f.\ $(39{,}876; 31{,}158)$, $p<0.001$). \\
z-test for two means & needs the population $\sigma$, which we do not know \\
t-test on the correlation of year built with \$/sq ft & sees only a straight-line trend over all years, not the pre-war contrast, and ignores location \\
\bottomrule
\end{tabularx}}

\tpart{c}{Result}
{\small
\begin{tabularx}{\linewidth}{@{}K{3.4cm} Y@{}}
\toprule
Estimate & +0.43\% (367 cells, 58,139 sales) \\
95\% confidence interval & $-1.68$\% to $+2.58$\% \\
Test statistic (d.f.) & $t=0.394$ (366) \\
p-value & 0.35 \\
Decision & do not reject $H_0$ \\
\bottomrule
\end{tabularx}}

Like for like, there is no pre-war premium: the whole interval lies inside $\pm5$\%. The 7.1\%
pooled gap is location and building mix, not age. Without year built 1935--1945 (the heaping at
1940) the estimate is +2.1\% (95\% CI $-0.3$\% to $+4.5$\%, $t=1.70$, d.f.\ 346, $p=0.045$):
borderline significant, still below what matters.

%------------------------------------------------------------ Test 2
\subsection*{Test 2 (P2): did the 2019 mansion tax push deals just below \$2M?}

\tpart{a}{Hypotheses, written before running the test}
{\small
\begin{tabularx}{\linewidth}{@{}K{3.4cm} Y@{}}
\toprule
$H_0$ & $p_{\text{post}}=p_{\text{pre}}$ \\
$H_1$ & $p_{\text{post}}>p_{\text{pre}}$ \\
Significance level & $\alpha=0.05$ \\
Side & one-sided (right) \\
Test & z-test for two proportions (pooled) \\
Smallest difference that matters & 10 percentage points \\
\bottomrule
\end{tabularx}}

$p$ = share of residential deals (tax class 1--2, one row per deal) within 2.5\% of \$2M that are
below \$2M; pre = 2016--18, post = 2020--25, 2019 dropped (contracts signed before April kept the
old rules). Since July 2019 crossing \$2M raises the tax on the whole price (about \$5,000), and
the tax only rewards being below, so the test is one-sided. 10 pp is about one buyer in ten moving
their price for the tax.

\tpart{b}{Other tests we considered}
Two independent groups of deals, a yes/no outcome and counts in the hundreds (all four cells
$\ge15$): a z-test for two proportions, with the rate pooled because $H_0$ says both periods share
one.
{\footnotesize
\begin{tabularx}{\linewidth}{@{}L{2.4cm} Y@{}}
\toprule
\textbf{Also considered} & \textbf{Why we did not use it} \\
\midrule
$\chi^2$ test of independence (2$\times$2) & same evidence but two-sided only. Run anyway: $\chi^2=54.85=z^2$, d.f.\ 1, $p<10^{-12}$. \\
z-test for one proportion vs 50\% & 50/50 is arbitrary: round prices at exactly \$2,000,000 make the pre-tax split uneven, so the pre-period is the right baseline \\
\bottomrule
\end{tabularx}}

\tpart{c}{Result}
{\small
\begin{tabularx}{\linewidth}{@{}K{3.4cm} Y@{}}
\toprule
Estimate & 47.0\% (157/334) $\rightarrow$ 69.6\% (673/967): +22.6 pp \\
95\% confidence interval & +16.5 to +28.7 pp \\
Test statistic & $z=7.41$ \\
p-value & $6.5\times10^{-14}$ \\
Decision & reject $H_0$ \\
\bottomrule
\end{tabularx}}

After the tax, deals pile up just below \$2M, and the whole interval is above the 10 pp that
matters. The placebos behave as they should: at \$1.5M (+1.7 pp, $p=0.25$) and \$2.5M (+0.5 pp,
$p=0.45$), where the tax did not change, nothing moved. At \$3M, the next tax step, the share rose
17.5 pp (95\% CI +9.8 to +25.3, $z=4.43$, $p<10^{-5}$).

%------------------------------------------------------------ Test 3
\subsection*{Test 3 (P3): do boroughs differ in \$0 transfers once building type is fixed?}

\tpart{a}{Hypotheses, written before running the test}
{\small
\begin{tabularx}{\linewidth}{@{}K{3.4cm} Y@{}}
\toprule
$H_0$ & $O_b=E_b$ for every borough \\
$H_1$ & at least one borough differs \\
Significance level & $\alpha=0.05$ \\
Side & upper tail of $\chi^2$ \\
Test & $\chi^2$ goodness of fit, d.f.\ $=5-1=4$ \\
Smallest difference that matters & 2 pp from a borough's expected share \\
\bottomrule
\end{tabularx}}

$O_b$ = \$0--\$100 transfers in borough $b$; $E_b=\sum_c r_c\,n_{bc}$, where $r_c$ is the citywide
transfer rate of building category $c$ and $n_{bc}$ the deeds of that category in borough $b$. So
$E_b$ is what the borough's building mix alone predicts.

\tpart{b}{Other tests we considered}
The data are counts in five categories compared with a predicted split: goodness of fit. Building
$E_b$ from category rates is what holds building type fixed.
{\footnotesize
\begin{tabularx}{\linewidth}{@{}L{2.4cm} Y@{}}
\toprule
\textbf{Also considered} & \textbf{Why we did not use it} \\
\midrule
$\chi^2$ test of independence, borough $\times$ transfer & mixes building type into the borough effect. Run anyway: $\chi^2=4{,}473$, d.f.\ 4, raw shares 21.1\% (Manhattan) to 37.7\% (Brooklyn), a 16.6 pp spread. \\
z-tests for two proportions, pair by pair & 10 pairs means 10 more tests and 10 more chances of a false alarm \\
$\chi^2$ test of independence inside each building category & holds type fixed too, but means 47 separate tests and many categories with tiny counts \\
\bottomrule
\end{tabularx}}

\tpart{c}{Result}
{\small
\begin{tabularx}{\linewidth}{@{}K{3.4cm} Y@{}}
\toprule
Test statistic (d.f.) & $\chi^2=461.3$ (4); smallest $E_b=7{,}873$ \\
p-value & $<10^{-90}$ \\
Decision & reject $H_0$ \\
Observed minus expected & Bronx $-4.9$ pp ($-5.5$ to $-4.3$), Staten Island $-5.5$ ($-6.1$ to $-4.9$), Brooklyn $+2.0$ ($+1.6$ to $+2.3$), Queens $+1.4$ ($+1.1$ to $+1.7$), Manhattan $-0.2$ ($-0.5$ to $+0.2$) \\
\bottomrule
\end{tabularx}}

Holding building type fixed changes the story. Manhattan's low raw share (21.1\%) is fully
explained by what sells there (condos and co-ops): it sits 0.2 pp from its expectation. The gaps
that matter are in the Bronx and Staten Island, with about 5 pp \emph{fewer} transfers than their
building mix predicts. Expected counts use rates from the same data, so d.f.\ $=4$ is approximate.

%------------------------------------------------------------ Test 4
\subsection*{Test 4 (P4): do buildings with shops sell at a different price per sq ft?}

\tpart{a}{Hypotheses, written before running the test}
{\small
\begin{tabularx}{\linewidth}{@{}K{3.4cm} Y@{}}
\toprule
$H_0$ & $\mu_d=0$ \\
$H_1$ & $\mu_d\ne0$ \\
Significance level & $\alpha=0.05$ \\
Side & two-sided \\
Test & paired t-test \\
Smallest difference that matters & 5\% either way \\
\bottomrule
\end{tabularx}}

$d$ = mean log \$/sq ft of buildings with commercial units minus purely residential ones, per
neighbourhood $\times$ type cell ($\ge3$ of each; mixed-use sales are rare). Two-sided: shops can
add rent (premium) or noise and traffic (discount).

\tpart{b}{Other tests we considered}
The same matched design as Test~1: one pair per cell, $\sigma$ unknown, 225 cells.
{\footnotesize
\begin{tabularx}{\linewidth}{@{}L{2.4cm} Y@{}}
\toprule
\textbf{Also considered} & \textbf{Why we did not use it} \\
\midrule
Welch two-sample t on all sales & mixed-use buildings sit on commercial streets in other neighbourhoods. Run anyway: $-9.3$\% ($-11.9$\% to $-6.7$\%), $t=-6.79$, d.f.\ 3,370. \\
Pooled two-sample t & the F-test rejects equal variances ($F=2.23$, $p<0.001$) \\
t-test on the correlation of commercial units with \$/sq ft & linear only, and most buildings have 0 or 1 commercial unit \\
\bottomrule
\end{tabularx}}

\tpart{c}{Result}
{\small
\begin{tabularx}{\linewidth}{@{}K{3.4cm} Y@{}}
\toprule
Estimate & $-15.5$\% (225 cells, 41,093 sales) \\
95\% confidence interval & $-19.5$\% to $-11.2$\% \\
Test statistic (d.f.) & $t=-6.74$ (224) \\
p-value & $1.3\times10^{-10}$ \\
Decision & reject $H_0$ \\
\bottomrule
\end{tabularx}}

Buildings with shops sell for about 15\% less per sq ft than purely residential buildings of the
same type in the same neighbourhood, and the whole interval is beyond 5\%. Ignoring neighbourhood
understates it ($-9.3$\%), because mixed-use buildings are more common in pricier areas.

%-----------------------------------------------------------------------
\section{Limitations}

\begin{itemize}
  \item \textbf{Association, not cause}, for P1, P3 and P4. In P4, gross floor area includes the
  shop space, which may simply be valued lower than homes. P2's placebos support a tax effect, but
  cannot rule out money moved off the deed.
  \item \textbf{The toll question was not tested.} Before the toll, the yearly gap between inside
  and outside the zone already swung from $+9.2$ to $-3.6$ pp (2021--24), so 2025's $-5.0$ pp cannot
  be told apart from past movement; a before/after test would wrongly credit the toll.
  \item \textbf{Price per sq ft covers houses and rental buildings only} (29.4\% of raw rows); condo
  and co-op units carry no unit area, so P1 and P4 say nothing about apartments.
  \item \textbf{Year built is an assessor estimate}, heaped on round years (Fig.~\ref{fig:heaping});
  dropping 1935--1945 moves T1 from +0.4\% to +2.1\%, so ``pre-war'' carries measurement error.
  \item \textbf{Sales are not independent.} Sales in one building, neighbourhood or month share
  conditions. Tests 1 and 4 work on cells to soften this; Tests 2 and 3 treat deals as independent,
  so their p-values are somewhat optimistic.
  \item \textbf{Judgement calls.} The \$100 and \$10,000 cut-offs, the multi-lot rule, the 2.5\%
  window and the minimum cell sizes were fixed before testing but are still choices.
  \item \textbf{Only properties that sold} are in the data; they may differ from the housing stock
  as a whole.
\end{itemize}

\end{document}
